Un organismo encargado de planificar programas de nutrición deportiva desea establecer un sistema de valoración económica de los requerimientos dietéticos de sus atletas.

Se consideran:

\begin{itemize}
    \item $m$ nutrientes esenciales indexados por $i$ ($i=1,\ldots,m$), para los que es necesario alcanzar al menos $b_i$ unidades.
    \item $p$ componentes restringidos indexados por $k$ ($k=1,\ldots,p$), cuya presencia en la dieta no debe superar $q_k$ unidades.
    \item $n$ alimentos indexados por $j$ ($j=1,\ldots,n$), cada uno con un coste unitario $c_j$.
\end{itemize}

Se introducen las siguientes variables de valoración:

\begin{itemize}
    \item $y_i$: valoración marginal asociada al cumplimiento del requerimiento mínimo del nutriente $i$.
    \item $w_k$: penalización marginal asociada a la restricción sobre el componente crítico $k$.
\end{itemize}

El objetivo es determinar dichas valoraciones de forma que el valor total neto de los requerimientos nutricionales sea máximo, respetando que la valoración imputada a cada alimento no exceda su coste real.

El modelo de programación lineal se formula como:

\[
(P)\qquad
\max \quad
z=
\sum_{i=1}^{m} b_i y_i
-
\sum_{k=1}^{p} q_k w_k
\]

sujeto a

\[
\sum_{i=1}^{m} a_{ij}y_i
-
\sum_{k=1}^{p} s_{kj}w_k
\le c_j,
\qquad
j=1,\ldots,n,
\]

\[
y_i \ge 0,
\qquad
i=1,\ldots,m,
\]

\[
w_k \ge 0,
\qquad
k=1,\ldots,p.
\]

\begin{enumerate}
    \item Formula el problema dual de $(P)$.
    \item Interpreta económicamente las variables del problema dual obtenido.
\end{enumerate}